\cotefile{23}
\scanpage{17}
% L1–3, vertical pen stroke in the left margin
En particulier, si $\Lambda$ est $k$-libre de rang fini, on obtient $\unsure{\mathfrak{g}|g|\mathcal{T}}_{A/k}\otimes_k I(\Lambda)$\inserted{$\simeq\pi_0(\mathfrak{g}\otimes_k U)$}\unsure{;|,} il en est de $=$ \deleted{\unsure{de}} si $\Omega^1_{A/k}$ est libre de rang fini sur $A$.

% L4
[alors $A/k$ est \unsure{diff$^{\rm t}$|diff.} \unsure{simplement|simple}] 
% archive sheet number: 16

\scanpage{18}
% L1–2, upper right
Fibrés des germes \deleted{\unsure{de|d}}d'ordre $n$ de \uline{sections}, etc.

% L3–4
Transformations \uline{infinitésimales}
% archive sheet number: 17

\scanpage{19}
% L1
$S$ préschéma

% L2–5
$\Lambda$ faisceau loc. libre de rang fini de $\underline{O}_S$-algèbres $\underline{O}_S$-augmentées, \deleted{tel que} d'idéal d'augmentation $I$, tel que $\operatorname{gr}^{I}(\Lambda)$ soit encore localement libre [donc $I^n=0$ pour $n$ \unsure{grand|grd}]. Un tel $\Lambda$ sera dit ``\uline{faisceau de type infinitésimal}''

% L6
$T=\operatorname{Spec}(\Lambda)$, $S$-schéma fini et plat, muni d'une section $S\to T$

% L7–8
\[
\begin{aligned}
\Lambda_n&=\Lambda/I^{n+1},\quad T_n=\operatorname{Spec}(\Lambda_n)\quad\text{d'où}\\
T_0&\simeq S\subset T_1\subset T_2\ldots\subset T_n\subset\ldots T_N^{\deleted{\unsure{n}}}=T
\end{aligned}
\]

% L7–9, at left, diagram enclosed on three sides
\diagram{\begin{tikzcd}[ampersand replacement=\&,execute at end picture={\draw (current bounding box.north west) -- (current bounding box.north east) -- (current bounding box.south east) -- (current bounding box.south west);}]
\& X\arrow[d,"f"'] \& X_0\arrow[l]\arrow[d,"f_0"]\\
S \& T\arrow[l] \& S=T_0\arrow[l]
\end{tikzcd}}
% slanted
\margin{\unsure{donnons nous au} \unsure{préalable} $X_0$}

% L9
$X$ un $T$-\inserted{pré}schéma, \inserted{\uline{simple}} sur $T$, \quad $X_n=X\times_T T_n\deleted{\unsure{X}}$

% L10–11
On considère de plus
\[
\xi=X_0\mathbin{\deleted{\otimes}\times}_S T=X_0\otimes_{\underline{O}_S}\Lambda
\]
\margin{Donc $\xi_n=X_0\times_S T_n$ en particulier $\xi_0=X_0$}
% slanted
\margin{mettre les hypothèses pour \unsure{simplicité|simplifier}, \unsure{publier} en \unsure{Annales|tête} \unsure{pour} \unsure{l'énoncé|l'idée} de 2 \uline{\unsure{propositions|présentations}}}

% L12–13
Du point de vue espaces topologiques, on a
\[
|X|=|\xi|=|X_0|=|\xi_0|
\]

% L14–16
$\underline{O}_X$ est un faisceau \inserted{de $f_0^\ast(\underline{\Lambda})$-algèbres} sur $X_0$, augmenté vers $\underline{O}_{X_0}$,
\diagram{\begin{tikzcd}[ampersand replacement=\&]
\&\underline{O}_X\arrow[r] \& \underline{O}_{X_0}\\
\underline{O}_S\arrow[r] \& \underline{\Lambda}\arrow[u]\arrow[r,"\varepsilon"] \& \underline{O}_S\arrow[u]
\end{tikzcd}\quad avec commutativité}

% L17–18
de plus dans \deleted{\unsure{f|la}} filtration \inserted{$I$-adique} de $\underline{O}_X$, \deleted{\unsure{pour}} le gradué associé est loc. libre sur $\underline{O}_{X_0}$ ($X$ plat sur $T$).

% L19–21
\deleted{\unsure{Soit|Si}} Pour tt ouvert $U\subset X_0=|X|$, soit $\operatorname{Aut}^{X_0}_{T}(U)$ ou $\operatorname{Aut}^{X_0}_{\underline{\Lambda}}(X)$ le groupe des $T$-automorphismes de $U$ induisant l'identité sur $X_0\cap U=U\times_T T_0$.
% archive sheet number: 18

\scanpage{20}
% L1–5
Ce groupe s'identifie aussi au groupe des $f_0^*(\underline{\Lambda})$-automorphismes de $\underline{O}_X\mid U$, qui induisent l'identité dans $\underline{O}_X/\underline{I}.\underline{O}_X$ ; donc quand $U$ varie, ces groupes forment les sections d'un faisceau de groupes, noté
% L6
\[
\underline{\operatorname{Aut}}^{X_0}_T(X)\quad\text{ou}\quad\underline{\operatorname{Aut}}^{X_0}_{\underline{\Lambda}}(X)
\]
% L7–9
qui s'identifie aussi au \deleted{faisceau} de \deleted{\unsure{des}} groupes \deleted{des} \inserted{d'automorphismes de la} $f_0^*(\underline{\Lambda})$-\deleted{Algèbre} $\underline{O}_{X_0}$-augmentée $\underline{O}_X$.

% L10
D'autre part, considérons \deleted{pour tout ouvert} \inserted{le $S$-schéma}
% L11
\[
\deleted{\unsure{X}}\qquad\overline X=\underline{\operatorname{Hom}}_{T/S}(T\deleted{\unsure{_0}},X)\deleted{\unsure{,}}
\]
\margin{``schéma des sections de $X$ sur $\Lambda$''}
% L12–15
La \deleted{\unsure{bonne|homom.}} section $S\xrightarrow{\sim}T_0\longrightarrow T$ définit
\[
\underline{\operatorname{Hom}}_{T/S}(T,X)\longrightarrow\underline{\operatorname{Hom}}_{T_0/S}(T_0,X\times_T T_0)\simeq X_0,
\]
d'où un homomorphisme de $S$-\inserted{pré}schémas
\[
\overline X\xrightarrow{\pi_\Lambda^X}X_0
\]
% L16–18
Si $U$ est un ouvert de $X_{\deleted{\unsure{0}}\unsure{0|{}}}$, alors on a
\[
\deleted{U}\qquad\overline U\simeq(\pi_\Lambda^X)^{-1}(U)
\]
(compatible avec \unsure{la projection|les projections}).
% L19–20
Par suite, $\underline{\operatorname{Aut}}^{X_0}_T(X)$ \uline{opère comme faisceau d'automorphismes sur le} $X$-\uline{schéma} $\overline X$
% archive sheet number: 19

\scanpage{21}
% L1–2
\uline{et par suite aussi sur le faisceau} $\boxed{\underline S(\overline X/X_0)}$ \uline{des germes de sections de} $\overline X$ \uline{sur} $X_0$

% L3–8
On peut aussi considérer le faisceau $\underline{\operatorname{Is}}_T^{X_0}(\deleted{X}\xi,\deleted{\xi}X)$, qui est un faisceau \uline{principal homogène} (\unsure{référence|v. finesse|références|vérifiant} \ldots) sous $\underline{\operatorname{Aut}}_T^{X_0}(X)$ (opérant à gauche) et sous $\underline{\operatorname{Aut}}_T^{X_0}(\xi)$ (\rule{2.5cm}{0.35pt} droite). D'ailleurs on a une section privilégiée \inserted{$s_0$} de $\underline S(\overline{\xi}/\xi_0)$

% L9–11
et on en déduit un hom de faisceaux
\[
u\longmapsto u(s_0)
\]
\[
\overline{\underline{\operatorname{Is}}}_T^{X_0}(\xi,X)\longrightarrow\underline S(\overline X/X_0)
\]

% L12–15, vertical rule to the left
\uline{Proposition} L'homomorphisme précédent est un isomorphisme de faisceaux, \deleted{\unsure{de groupes à}} \inserted{compatible avec le} faisceau de groupes d'opérateurs $\underline{\operatorname{Aut}}_T^{X_0}(X)$ (à gauche).

% L16–17
\uline{Corollaire 1} $\underline S(\overline X/X_0)$ est \uline{principal homogène} sous le faisceau $\underline{\operatorname{Aut}}_T^{X_0}(X)$

% L17–19, left
% slanted
\margin{Ceci vaut $=$ si $\underline{\Lambda}$ \unsure{---} $\longleftrightarrow$ \unsure{qu'une|par une} $\underline{O}_S$-Algèbre, \unsure{mais} \unsure{---} sans $\underline{O}_S$-augmentation, cf. Greenberg pour la définition de $\overline X$}
% L18–20
\uline{Corollaire 2} $\underline{\operatorname{Aut}}_T^{X_0}(\xi)$ opère de façon naturelle \inserted{à droite} sur $\underline S(\overline X/X_0)$, qui devient un \uline{faisceau principal hom.} à droite sous ce faisceau de groupes \deleted{\unsure{à droite}}
% archive sheet number: 20

\scanpage{22}
% L1–2
\uline{Corollaire 3} On a un hom naturel de faisceaux
\[
\underline S(\overline X/X_0)\longrightarrow\underline{\operatorname{Is}}_{X_0}(\overline{\deleted{X}\xi},\overline X)
\]
% L3–7
compatible avec les opérations de $\underline{\operatorname{Aut}}_T^{X_0}(X)$ à gauche, et les opérations \unsure{---|à d.} de $\underline{\operatorname{Aut}}_T^{X_0}(\xi)$ [qui opère sur le premier \unsure{mb|m$^{\rm bre}$} par le corollaire 2]. Cet homomorphisme est \uline{injectif} [car $\underline{\operatorname{Aut}}_T^{X_0}(\xi)$ opère fidèlement dans $\overline\xi$]

% L8–12
On peut dire aussi [mais cela n'est pas très \unsure{malin|utile}] que $\overline X/X_0$ \uline{est muni d'une structure de} \deleted{fais} \uline{schéma sur} $X_0$ \uline{à faisceau structural} $\unsure{=|:}\boxed{\underline{\operatorname{Aut}}_T^{X_0}(\xi)=\underline G_{X_0}^{\underline{\Lambda}}}$ \uline{opérant sur le schéma-type} $\xi=X_0\times_S T=X_0\otimes_{\underline{O}_S}\underline{\Lambda}$

% L13–16
\uline{La donnée de} $X$ \uline{simple sur} $T$, se \uline{réduisant suivant} $X_0$, est \uuline{équivalente} \uline{à la donnée d'un schéma} $\overline X$ \uline{sur} $X_0$, \uline{muni d'un tel faisceau structural}.

% L17–18
\uline{Le faisceau principal associé n'est autre que} $\underline S(\overline X/X_0)$ !

% L19–21
$\overline\xi$ est noté aussi $X_0^{\{\Lambda\}}$, et appelé \deleted{\unsure{S-schéma}} \uline{$S$-préschéma des p$^{\rm ts}$ proches} de $X_0$ \uline{du type} $\Lambda$. \unsure{Ses|les} sections \deleted{\unsure{sont des}} forment donc
% archive sheet number: 21

\scanpage{23}
% L1–4
s'interprètent, \uline{soit} comme un $\unsure{\Lambda|T|{}}$-automorphisme de $X_0\mathbin{\deleted{\otimes}\times}_S T=X_0\otimes_{\underline{O}_S}\underline\Lambda$ induisant l'identité sur $X_0$, soit comme la donnée d'un champ de sections \deleted{\unsure{proches|privilégiées} $\Lambda$ \illegible{2}} \ldots\ldots
% L5, detached stroke at left, middle third
\margin{\unsure{---}}
% archive sheet number: 22

\scanpage{24}
% L1, light pencil
\margin{$\deleted{\unsure{(X/Y)^p}}=\unsure{(X/Y)^p}=\infty\qquad\unsure{Z|2}$}

% L1–3, upper left
\diagram{\begin{tikzcd}[ampersand replacement=\&]
\&Y\arrow[d,"p"]\&\pi^{(n)}(Y/X/S)\arrow[d]\\
S\&X\arrow[l,"f"]\&\Delta^{(n)}_{X/S}\arrow[l,"{\unsure{g|p|D|\zeta|\delta}}^{n}_{X/S}"]
\end{tikzcd}}
% L3–5
On veut définir les ``fibrés''
\[
p^n(Y/X/S)
\]
des \deleted{éléments} ``éléments de sections d'ordre $n$ de $Y/S$ \unsure{en|sur} $X/S$''.

% L5–9, left
\diagram{\begin{tikzcd}[ampersand replacement=\&]
\&X\times_S Y\arrow[d]\\
X\arrow[r,"\Delta_X"]\&X\times_S X\arrow[d,"\operatorname{pr}_1"]\\
\&X
\end{tikzcd}}
% L6–9
\inserted{$X\times_S X$ est} muni des deux projections $\operatorname{pr}^{\deleted{\unsure{2}}}_1$, $\operatorname{pr}_2$ et de l'application diagonale $\Delta_X$. On considère $\deleted{\unsure{\delta}}\ \Delta^{(n)}_{X/S}$ le $n$-ième voisinage infinitésimal de l'immersion $\Delta_X$ comme un $X$-\deleted{\unsure{p|pré}}schéma grâce à $\operatorname{pr}_1$.

% L9–12, left
\diagram{\begin{tikzcd}[ampersand replacement=\&]M\arrow[d]\\
{\unsure{P}}^{n}_{A/\Lambda}\otimes_A M\arrow[d]\\ M
\end{tikzcd}}
% L10–13
$i^{(n)}_X$ désigne l'immersion de $\Delta^{(n)}_{X/S}$ \deleted{dans} $X\times_S X$. On considère alors le
\[
\deleted{p^n(Y/X/S)=(\operatorname{pr}_2 i^{(n)}_X)^\ast}
\]
$\Delta^{(n)}_{X/S}$-schéma $(\operatorname{pr}_2 i^{(n)}_X)^\ast(Y/X)\inserted{=\Delta^n_{X/S}\times_{X\times_S X}X\times_S Y}=\deleted{\unsure{p}}\pi^{(n)}(Y/X/S)$

% L14–17, diagram at left
\diagram{\begin{tikzcd}[ampersand replacement=\&,column sep=small,row sep=small,execute at end picture={\draw ([xshift=-4pt,yshift=4pt]\tikzcdmatrixname-1-3.north west) rectangle ([xshift=4pt,yshift=-5pt]\tikzcdmatrixname-2-4.south east);}]
\vphantom X\&X\times_S Y\arrow[l]\arrow[d]\&p^n(Y/X/S)\arrow[l]\arrow[d]\&Y\arrow[l,"{\varepsilon}^{n}(Y/X/S)"']\arrow[d]\\
X\&X\times_S X\arrow[l,"\operatorname{pr}_1"']\&\Delta^{(n)}_{X/S}\arrow[l,"i^{(n)}_{X/S}"']\&X\arrow[l]\\
X'\arrow[u,"\varphi"]\&X'\times_S X\arrow[u]\arrow[l,"\unsure{\mathrm{id}|i\Lambda}"']\&\&X'\arrow[u]\arrow[ll,"(\Gamma_{\varphi})^{(n)}"']
\end{tikzcd}}
% L14–16
On a donc $X\times_{\Delta^{(n)}_{X/S}}\pi^n(Y/X/S)=Y$, d'où $Y\longrightarrow\pi^n(Y/X/S)$ compatible avec
\[
X\xrightarrow{\Delta_X}\Delta^{(n)}_{X/S}
\]
% L17–19
Il y a aussi un \deleted{le} morphisme composé
\[
d^{(n)}_{Y/X/S}:\pi^n(Y/X/S)\longrightarrow X\times_S Y\longrightarrow Y
\]
et on a commutativité de
% L20–21
\diagram{\begin{tikzcd}[ampersand replacement=\&,execute at end picture={\draw ([xshift=-5pt,yshift=5pt]\tikzcdmatrixname-1-2.north west) rectangle ([xshift=5pt,yshift=-5pt]\tikzcdmatrixname-2-3.south east);}]
Y\arrow[r,"\varepsilon^n(Y/X/S)"]\arrow[d]\&\pi^n(Y/X/S)\arrow[d]\arrow[r,"d^n(Y/X/S)"]\&Y\arrow[d]\\
X\arrow[r,"\Delta_{X/S}"]\&\Delta^{(n)}_{X/S}\arrow[r,"d^{(n)}_{X/S}"]\&\deleted{\unsure{\Delta^n_{X/S}}}X
\end{tikzcd}}
% archive sheet number: 23

\scanpage{25}
% L1
Ce n'est donc pas un $\deleted{\unsure{S|X}}X$-morphisme en général.

% L2–4
Supposons maintenant que $\Delta^{(n)}_{X/S}$ soit \deleted{\unsure{plat fini sur|plat, de présentation}} \inserted{fini} \deleted{\unsure{et}} loc. libre sur $X$. Alors on peut former
\[
\boxed{({\unsure{g|D|\zeta|\delta}}^n_{X/S})_*(\pi^n(Y/X/S)/\Delta^{(n)}_{X/S})=P^{(n)}(Y/X/S).}
\]
% slanted
\margin{\uline{fibré} \uline{des éléments} \uline{de sections d'ordre $n$} \uline{de} $Y/S$ \uline{sur} $X/S$}

% L5–6
Par définition, on aura donc
\[
\Gamma(\pi^n(Y/X/S)/\Delta^{(n)}_{X/S})\simeq\Gamma(p^{(n)}(Y/X/S)/X)
\]
% L7–9
et \deleted{\unsure{le foncteur|de façon}} plus généralement, si \inserted{$\varphi:$} $X'\longrightarrow X$ un préschéma sur $X$
\[
\begin{aligned}
\operatorname{Hom}_X(X',P^{(n)}(Y/X/S))&\simeq\operatorname{Hom}_{\Delta^{(n)}_{X/S}}(X'\times_X\Delta^n_{X/S},\pi^{(n)}(Y/X/S)\\
&\simeq
\end{aligned}
\]
% L10–12
Un \unsure{él$^{\rm mt}$|él$^{\rm nt}$} de cet ensemble est donc aussi obtenu en prenant l'immersion
\[
\Gamma_\varphi=\deleted{\unsure{\operatorname{id}}}\varphi\times_S\deleted{\unsure{\varphi}}\operatorname{id}_X:X'\longrightarrow X'\times_S X
\]
\margin{$\deleted{\unsure{X\times_S X}}$}
% L13–16
en prenant son $n$.ième \unsure{voisinage} d'immersion
\[
\Gamma_\varphi^{(n)}\longrightarrow X'\times_S X
\]
et une section \uline{de $X'$}$\times_S Y$ là-dessus ///

% L17–19
On \deleted{\unsure{va donc|va alors}}

On a un $X$-morphisme canonique \unsure{encore} noté $D^n(Y/X/S)$ :
% archive sheet number: 24

\scanpage{26}
% L1–2
\diagram{\begin{tikzcd}[ampersand replacement=\&]
p^{(n)}(Y/X/S)\arrow[rr,"{\unsure{\varepsilon|E}}^n(Y/X/S)"]\arrow[dr]\&\&Y\arrow[dl]\\
\&X\&
\end{tikzcd}}
% L2–4
d'autre part, on va définir un \uline{hom} \deleted{application} \uline{de faisceaux} sur $X$
\[
d^n(Y/X/S):\underline{\Gamma}(Y/X)\longrightarrow\underline{\Gamma}(P^{(n)}(Y/X/S)/X)
\]
% slanted
\margin{\uline{\deleted{\unsure{système de}}} \uline{parties principales}}
% L5–8
[tel que le composé avec $\Gamma$
\[
{\unsure{g|D|\delta|\zeta|\varepsilon}}^n(Y/X/S):\underline{\Gamma}(p^{(n)}(Y/X/S)/X)\longrightarrow\underline{\Gamma}(Y/X)
\]
soit l'identité :
\[
g^n d^n=\operatorname{id}\qquad\big]
\]

% L8, small isolated mark below display, at left
\margin{$\unsure{x|\times}$}
% L9–12
\deleted{Il faut donc}

Pour ceci, \deleted{\unsure{on note que}} \inserted{il suffit de définir}
\[
\Gamma(Y/X)\longrightarrow\Gamma(P^n(Y/X/S)/X)
\]
et de vérifier sa nature locale sur $X$. On note que
% L13–14
\[
\Gamma(Y/X)\longrightarrow\deleted{\operatorname{Hom}}\Gamma(Y\times_X\Delta^{(n)}_{X/S}/\Delta^{(n)}_{X/S})
\]
% L15–20
Or
\[
\begin{aligned}
\operatorname{Hom}_X(\Delta^n_{X/S},Y)&=\Gamma(Y\times_X\Delta^{(n)}_{X/S}/\Delta^{(n)}_{X/S})\simeq\text{\unsure{schéma|ensemble|réunion} des}\\
&\text{sections d'\unsure{ordre|ord.|au-dessus}\deleted{\unsure{x}} $n$ de }Y\times_S X\text{ sur}\\
&X\times_S X\text{ le long de l'immersion}\\
&X\longrightarrow X\times_S X
\end{aligned}
\]
et par \unsure{symétrie|hypothèse|projections}, ce dernier ensemble est
% archive sheet number: 25

\scanpage{27}
% L1, cancelled formula at left
\margin{$\deleted{M\xrightarrow{d^n}P^n_{A/\Lambda}(M)}$}
% L1–3
en correspondance biunivoque avec les éléments de \unsure{section|sections} d'ordre $n$ de $X\times_S Y$ \deleted{le long} sur $X\times_S X$ le long de la $=$ immersion \unsure{\ldots|{}}. \deleted{\unsure{$\Gamma_\varphi$}}

% L4–8
Soit $n'\geq n$, alors on a
\[
\Delta^{(\deleted{\unsure{n'|m}}n)}_{X/S}\longrightarrow\Delta^{(n')}_{X/S}
\]
et les $\Delta^{(n)}_{X/S}$ forment un \uline{système inductif} de schémas sur $X$, dont la limite est un \uline{schéma formel} sur $X^{\unsure{\to|{}}}$, noté $\Delta^{(\infty)}_{X/S}$.

% L9–12
On a évidemment
\[
\pi^n(Y/X/S)=\pi^{n'}(Y/X/S)\times_{\Delta^{(n')}_{X/S}}\Delta^{(n)}_{X/S}
\]
d'où aussitôt un $X$-hom naturel
\[
P^{n'}(Y/X/S)\xrightarrow{\varphi_{nn'}}p^n(Y/X/S).
\]
% L13–16
C'est un système \uline{transitif} de $X$-hom. et on a \deleted{$\unsure{\varphi}$}
\[
\left\{\begin{aligned}
\underline{\varphi_{\unsure{0|\circ|o}n}=\varepsilon^n}&\quad\text{d'où}\\
\underline{\underline{\varepsilon^{(n)}}}(Y/X/S)\circ\varphi_{nn'}&=\varepsilon^{(n')}\\
\varphi_{nn'}d^{(n')}&=d^{(n)}.
\end{aligned}\right.
\]
% archive sheet number: 26

\scanpage{28}
% L1–2
Les sections de $p^n(\deleted{\unsure{X/}}Y/X/S)$ sur $Y$ sont des $X$-morphismes, donc

% L3–7, main diagram
\diagram{\small\begin{tikzcd}[ampersand replacement=\&,column sep=small,row sep=small]
Y^{\{\Delta^{(n)}_{X/S}\}}\arrow[dr,"\unsure{n|{}}"'] \& \& \& \&\\
p^n(Y/X/S)\arrow[r] \& Y\arrow[d,"p"]\arrow[ul,dashed,bend right=13] \& Y\times_S X\arrow[l]\arrow[d] \& \pi^{(n)}(Y/X/S)\arrow[l]\arrow[d] \& Y\arrow[l]\arrow[d]\\
\&X\&X\times_S X\arrow[l,"\operatorname{pr}_1"]\&\Delta^{(n)}_{X/S}\arrow[l]\&X\arrow[l]\arrow[ll,dashed,bend left=15,"\deleted{\unsure{\Delta_X}}\ \Delta_{X/S}"]\\
\&Y'\arrow[u,start anchor=north west,end anchor=south east]\&Y'\times_S X\arrow[l]\arrow[u,start anchor=north west,end anchor=south east]\&\Gamma^{(n)}_{\unsure{p|\varphi}}\arrow[l]\arrow[u]\&Y'\arrow[l,"\deleted{\unsure{Y}}" description]\arrow[u,"p"']\arrow[uu,bend right=24,"\operatorname{id}"']\arrow[ll,dashed,bend left=15,"\Gamma_{\unsure{X|p}/Y/S}"]
\end{tikzcd}}

% L8–9
\diagram{\begin{tikzcd}[ampersand replacement=\&]
\&\pi^n(Y/X/S)\arrow[d]\&Y\arrow[l]\arrow[d]\\
\deleted{\unsure{Y}}\arrow[r]\&Y^{\Delta^{(n)}_{X/S}}\&Y\arrow[l]
\end{tikzcd}}

% L10–13, left
\diagram{\begin{tikzcd}[ampersand replacement=\&,column sep=small]
\&\&B\&\\
\unsure{\Lambda|A}\arrow[r]\&\Big\Vert\& A\arrow[u]\arrow[r,bend left=12]\arrow[r,bend right=12]\arrow[dr]\& A\otimes A\,\Big|\\
\&\&\&A'
\end{tikzcd}}
% L10–12, right
\diagram{\begin{tikzcd}[ampersand replacement=\&]
\&A'\otimes B\\
(B\otimes_A A')\otimes A\arrow[d,leftrightarrow]\&\\
J\subset \deleted{\unsure{B}}A'\otimes A\arrow[r]\arrow[uur,bend right=12]\&\deleted{\unsure{B}}A'
\end{tikzcd}}
% L13–14, lower edge
\diagram{\begin{tikzcd}[ampersand replacement=\&]
(B\otimes_A\deleted{\unsure{A}}B)\otimes A \& B\otimes B\\
B\otimes A\arrow[u]\arrow[ur]\arrow[r]\&B
\end{tikzcd}}
% archive sheet number: 27

\scanpage{29}
% L1, upper left
\[
m\longrightarrow1\otimes m
\]
% L2–6, left
\unsure{de}
\diagram{\begin{tikzcd}[ampersand replacement=\&]
Y^{\Delta^{\unsure{n|*}}_{X/S}}\arrow[dr]\&\&\\
p^n(Y/X/S)\arrow[r]\&Y\arrow[d]\&\pi^n(Y/X/S)\arrow[l]\arrow[d]\\
\&X\&\Delta^{\unsure{n|*}}_{X/S}\arrow[l]
\end{tikzcd}}
% L1–7, right
$\Gamma(Y^{\{\Delta^n_{X/S}\}}/Y)$ \uline{groupe} des automorphismes infinitésimaux d'ordre $n$ du fibré $Y/X$
% L5–7, small interline
$=$ \unsure{grp|groupe} des \unsure{automorphismes|endomorphismes} de $\pi^n(Y/X/S)/\Delta^{\unsure{n|2}}_{X/S}$ qui ind$^{\rm t}$ l'identité sur $Y$

% L6–8
$\Gamma(p^n(Y/X/S)/Y)$ \uline{ensemble} des connections d'ordre $n$ du fibré $Y/X/S$

% L9–12
Si $Y/X$ est \uline{simple}, alors \deleted{l'ensemble des connections est} $\underline{\Gamma}(p^n(Y/X/S)/Y)$ est principal homogène \unsure{sous} $\underline{\Gamma}(Y^{\Delta^n_{X/S}}/Y)={\unsure{\operatorname{Aut}|\operatorname{Tr}}}^{(n)}(Y/X/S)$

% L13–14, left
\diagram{\begin{tikzcd}[ampersand replacement=\&]
p^n(Y/S)\&p^n(X/S)\arrow[l]\\
O_Y\arrow[u]\&O_X\arrow[l]\arrow[u]
\end{tikzcd}}
% L13–15, right
hom de $\underline{O}_{\deleted{\unsure{X}}Y}$-Algèbres augmentées
\[
\underline{\underline{p^n(Y/X)\longrightarrow f^*(p^n(X/S))}}
\]

% L16–18, left
\diagram{\begin{tikzcd}[ampersand replacement=\&,column sep=small]
f^\ast p^n(Y/S)\arrow[d]\&p^n(Y/S)\arrow[l,"?"']\arrow[d]\&f^\ast p^n(X/S)\arrow[l,Rightarrow]\arrow[d]\\
\underline{O}_X\&\underline{O}_X\arrow[l,"\unsure{\text{comment}|\text{commut}}?"']\&\underline{O}_X\arrow[l]
\end{tikzcd}}
% L16–19, right
hom de de $f^*(p^n(X/S))$-algèbres $\underline{O}_X$-augmentées
\[
p^n(Y/S)\unsure{\longrightarrow|-|\longleftarrow}f^*p^n(X/S)
\]
(pas nécessairement augmentées)
% archive sheet number: 28

\scanpage{30}
% L1–3, sheet cancelled with long diagonal strokes
\[
\deleted{p^n_{A/\Lambda}\otimes_A M\xrightarrow{\overset{\varphi}{\sim}}M\otimes_A p^n_{A/\Lambda}}\qquad
\deleted{\varphi(g.\underbrace{d^nf\otimes m}_{\unsure{g[{\otimes}fm]|g\otimes f}})=m\otimes f d^n g}
\]
\[
\deleted{\unsure{(}p^n_{A/\Lambda})_g\quad(P^n_{A/\Lambda})_d}\qquad
\deleted{\varphi(gd^n(fm))=fm\otimes}
\]
% L4–6
\[
\begin{aligned}
&\deleted{(A\otimes_\Lambda A)\otimes_A M\xrightarrow{\sim}M\otimes_A(A\otimes_\Lambda A)}\\
&\deleted{A\otimes_\Lambda M\qquad M\otimes_\Lambda A}
\end{aligned}
\]
% circled at left
\[
\deleted{\mathord{\tikz[baseline=(p.base)]{\node[inner sep=0pt] (p) {$p^n_{A/\Lambda}\otimes_A M$};\draw (p.center) ellipse (1.15cm and .4cm);}}}
\]

% L7–10, left
\[
\deleted{M\longrightarrow A\otimes_\Lambda M}\quad\deleted{\unsure{\text{et}}}
\]
\diagram{\deleted{\mbox{\begin{tikzcd}[ampersand replacement=\&]
Y\arrow[d,dash]\arrow[r,dashed,bend left=18,"?"]\&P(Y)\arrow[l,"\deleted{\unsure{\mathrm{id}}}"]\arrow[l,bend left=37]\\
X\&
\end{tikzcd}}}}
% L7–10, centre
\[
\deleted{\unsure{=}}\qquad\deleted{S_A(\underline{\underline{\underline A}}\otimes_\Lambda M)}\quad\deleted{\unsure{\equiv}}
\]
\diagram{\deleted{\mbox{\begin{tikzcd}[ampersand replacement=\&]
S_A(M)\arrow[r]\&\underline A\otimes_\Lambda S_\Lambda(M)\\
\deleted{\unsure{M}}\arrow[r]\arrow[ur]\&\deleted{\unsure{B|A}}A\otimes_\Lambda\underline{S_A(M)}
\end{tikzcd}}}}
% L8–10, right
\diagram{\deleted{\mbox{\begin{tikzcd}[ampersand replacement=\&]
\deleted{\unsure{\xi}}\&M\arrow[d,dash]\\
A\arrow[r,dash]\&B
\end{tikzcd}}}}
\[
\deleted{\operatorname{Hom}_{B}(M,B)}
\]

% L11
\[
\deleted{M\longrightarrow A\otimes_{\unsure{A|\Lambda}} M}
\]
% L12–14, left
\diagram{\deleted{\mbox{\begin{tikzcd}[ampersand replacement=\&]
Y=\operatorname{Sp}\bigl(S_{\unsure{A|X}}(M)\bigr)\arrow[d]\\
X=\operatorname{Spec}(A)
\end{tikzcd}}}}
% L12–14, centre
\[
\begin{aligned}
&\deleted{(A\otimes_A A)\otimes_A M=A\otimes_\Lambda M}\\
&\deleted{\Delta_X^{\ast}=\operatorname{Sp}(A\otimes_\Lambda\Lambda)}
\end{aligned}
\]
% L12–13, right
\[
\deleted{\underline A\otimes_\Lambda M\leadsto M}\qquad\deleted{\unsure{\check V}.}
\]
% archive sheet number: 29

\scanpage{31}
% L1–4
Soit $f:\deleted{\unsure{Y}}X\longrightarrow\deleted{\unsure{X}}Y$ un $S$-morphisme. On appelle \uline{transformation infinitésimale de type} $S'$ \uline{de} $f$ un couple $u,v$ de $S'$-automorphismes de $X_{S'},Y_{S'}$ \inserted{induisant l'identité sur $X,Y$} rendant commutatif le diagramme
% L5–6
\diagram{\begin{tikzcd}[ampersand replacement=\&]
X_{S'}\arrow[r,"u"]\arrow[d,"f_{S'}"']\&X_{S'}\arrow[d,"f_{S'}"]\\
Y_{S'}\arrow[r,"v"]\&Y_{S'}
\end{tikzcd}}
% L7–9
Les transformations infinitésimales \inserted{$w=(u,v)$} de type $S'$ de $f$ forment un \uline{groupe}, contenu dans le produit des groupes des autom. inf. de type $S'$ de $X$ et $Y$.

% L10–13
Soit $w=(u,v)$ une transformation infinitésimale de type $S'$ de $f$, et soit $g$ une section de $X/Y$. Considérons alors
\[
\widetilde w(g)=u\,g_{S'}\,v^{-1}
\]
% L14–20
c'est une section de $X_{S'}$ sur $Y_{S'}$ ``proche'' de $g$, donc correspond à \unsure{une} \deleted{\unsure{S-flèche|flèche|système} $\deleted{\unsure{X}}Y\times_S S'{\to}X$ \unsure{la} \unsure{flèche} tel que son composé avec $X\unsure{-|\to}Y$ soit $\operatorname{pr}_1$, i.e. un $S$-hom $Y\to X^{\{S'\}}$ tel que son composé avec $X^{\{S'\}}\unsure{-|\to}Y^{\{S'\}}$ soit la section \unsure{évidente} $Y\to Y^{\{S'\}}$ \unsure{section de}} élément de $\Gamma(X_Y^{\{Y_{S'}\}}/Y)$ ``proche'' de $g$.

% L21
Propriétés fonctorielles triviales \ldots\ldots
% archive sheet number: 30

\scanpage{32}
% L1–5, upper paragraph
1°) Le groupe des transformations infinitésimales \inserted{de type $S'/S$} de $X_S/Y_S$ \uline{opère} sur $\Gamma(X_{S'}/Y_{S'})$ des sections \unsure{proches} de type $S'$ des sections de $X/Y$, en induisant l'identité sur le quotient $\Gamma(X/Y)$ de ce dernier ensemble, de plus

% L6–10, formula with gloss running below the brace
\[
\Gamma(X/Y)\longrightarrow
\Gamma\underbrace{(X_{S'}/Y_{S'})}_{\substack{\text{ensemble}\\
\text{: groupe}\\ \text{d'opérateurs}\\
\operatorname{Aut}^{\{\Lambda\}}_{S}(X_{S'},Y_{S'})}}
\simeq\Gamma(X_Y^{\{Y_{S'}\}}/Y)
\]

% L11–17, central note
\uline{N.B.} $S,S'$ n'interviennent pas effectivement \deleted{\unsure{mais}} \inserted{[\unsure{seult|seulement} par} $X_{S'}\to Y_{S'}$] dans la définition \deleted{\unsure{de}} la première ligne, mais seulement dans la définition du groupe d'opérateurs sur $\Gamma(X_{S'}/Y_{S'})$, : les transformations inf. de type $S'$ de $X_{\deleted{\unsure{S}}S}/Y_S$ i.e. les $S'$-automorphismes de $Y_{S'}$ et $X_{S'}$ .....

% L18–20, lower paragraph
2°) Fixons un $w$. Un $Y$-hom $\displaystyle X\xrightarrow{h}X'$ est compatible avec des connexions $w=(u,v)$, $w'=(u',v)$ si on a commutativité \unsure{du|de}

% L21–22, diagram, and incomplete continuation at right
\diagram{\begin{tikzcd}[ampersand replacement=\&]
X_{S'}\arrow[r,"h_{S'}"]\arrow[d,"u"'] \& X'_{S'}\arrow[d,"u'"]\\
X_{S'}\arrow[r,"h_{S'}"] \& X'_{S'}
\end{tikzcd}}
\hfill i.e. $\deleted{\unsure{X'|h'}}$ \unsure{l}

% archive sheet number: 31

